\documentclass[xcolor=table]{beamer}
\usepackage{beamerthemesplit}
\usepackage{wrapfig}
\usetheme{SPbGU}
\usepackage{pdfpages}
\usepackage{amsmath}
\usepackage{cmap}
\usepackage[T2A]{fontenc}
\usepackage[utf8]{inputenc}
\usepackage[english]{babel}
\usepackage{indentfirst}
\usepackage{amsmath}
\usepackage{tikz}
\usepackage{multirow}
\usepackage[noend]{algpseudocode}
\usepackage{algorithm}
\usepackage{algorithmicx}
\usepackage{fancyvrb}
\usetikzlibrary{shapes,arrows}
\usepackage{tikz}
\usetikzlibrary{automata,positioning}
\usepackage{listings}
\usepackage{multicol}
\setbeamertemplate{itemize items}[circle]

\lstdefinelanguage{scala}{
  morekeywords={%
          abstract,case,catch,class,def,do,else,extends,%
          false,final,finally,for,forSome,if,implicit,import,lazy,%
          match,new,null,object,override,package,private,protected,%
          return,sealed,super,this,throw,trait,true,try,type,%
          val,var,while,with,yield},
  otherkeywords={=>,<-,<\%,<:,>:,\#,@},
  sensitive=true,
  morecomment=[l]{//},
  morecomment=[n]{/*}{*/},
  morestring=[b]",
  morestring=[b]',
  morestring=[b]"""
}[keywords,comments,strings]

\lstset{language=Scala}

\beamertemplatenavigationsymbolsempty

\title[Поиск путей с КС ограничениями]{Использование алгоритмов синтаксического анализа для поиска путей с контекстно-свободными ограничениями}
\institute[СПбГУ]{
JetBrains Research, Лаборатория языковых инструментов  \\
Санкт-Петербургский государственный университет
}

\author[Екатерина Вербицкая]{Екатерина Вербицкая \\
\and
  {\bfseries Научный руководитель:} д.т.н. Д.В. Кознов  \\
\and
  {\bfseries Рецензент:}  д.т.н. Ф.А. Новиков
}



\date{18 июня 2019}

\definecolor{orange}{RGB}{179,36,31}

\begin{document}
{
\begin{frame}[fragile]
  \begin{tabular}{p{5.5cm} p{5.5cm}}
   \begin{center}
      \includegraphics[height=1.5cm]{pictures/jetbrainsResearch.pdf}
    \end{center}
    &
    \begin{center}
      \includegraphics[height=1.5cm]{pictures/SPbGU_Logo.png}
    \end{center}
  \end{tabular}
  \titlepage
\end{frame}
}

\begin{frame}[fragile]
  \transwipe[direction=90]
  \frametitle{Поиск путей в графах}
\begin{minipage}[m]{0.42\linewidth}
\begin{tikzpicture}[shorten >=1pt,node distance=2cm,on grid,auto]
   \node[state] (a) {\textbf{X}};
   \node[state] (b) [left=of a] {B};
   \node[state] (c) [right=of a] {C};
   \node[state] (d) [below=of c] {D};
   \node[state] (e) [below=of b] {E};
   \node[state] (f) [below=of a] {F};
   \node[state] (g) [below=of f] {\textbf{Y}};
   \node[state] (h) [below=of d] {H};
    \path[->]
    (a) edge node {} (b)
    (a) edge node {} (c)
    (a) edge node {} (f)
    (b) edge node {} (e)
    (c) edge node {} (d)
    (d) edge node {} (h)
    (d) edge node {} (a)
    (e) edge node {} (f)
    (e) edge node {} (g)
    (f) edge node {} (b)
    (f) edge node {} (g)
    (h) edge node {} (g);
\end{tikzpicture}
\end{minipage}\hfill
\begin{minipage}[m]{0.5\linewidth}
\pause    Найти путь из  \textbf{X} в \textbf{Y}
\end{minipage}
\end{frame}

\begin{frame}[noframenumbering]
  \transwipe[direction=90]
  \frametitle{Поиск путей в графах}
\begin{minipage}[m]{0.42\linewidth}
\begin{tikzpicture}[shorten >=1pt,node distance=2cm,on grid,auto]
   \node[state] (a) {\textbf{X}};
   \node[state] (b) [left=of a] {B};
   \node[state] (c) [right=of a] {C};
   \node[state] (d) [below=of c] {D};
   \node[state] (e) [below=of b] {E};
   \node[state] (f) [below=of a] {F};
   \node[state] (g) [below=of f] {\textbf{Y}};
   \node[state] (h) [below=of d] {H};
    \path[->]
    (a) edge[red, line width=2pt] node {} (b)
    (a) edge node {} (c)
    (a) edge node {} (f)
    (b) edge[red, line width=2pt] node {} (e)
    (c) edge node {} (d)
    (d) edge node {} (h)
    (d) edge node {} (a)
    (e) edge[red, line width=2pt] node {} (f)
    (e) edge node {} (g)
    (f) edge node {} (b)
    (f) edge[red, line width=2pt] node {} (g)
    (h) edge node {} (g);
\end{tikzpicture}
\end{minipage}\hfill
\begin{minipage}[m]{0.5\linewidth}
Найти путь из  \textbf{X} в \textbf{Y}
\end{minipage}
\end{frame}

\begin{frame}[noframenumbering]
  \transwipe[direction=90]
  \frametitle{Поиск путей в графах}
\begin{minipage}[m]{0.42\linewidth}
\begin{tikzpicture}[shorten >=1pt,node distance=2cm,on grid,auto]
   \node[state] (a) {\textbf{X}};
   \node[state] (b) [left=of a] {B};
   \node[state] (c) [right=of a] {C};
   \node[state] (d) [below=of c] {D};
   \node[state] (e) [below=of b] {E};
   \node[state] (f) [below=of a] {F};
   \node[state] (g) [below=of f] {\textbf{Y}};
   \node[state] (h) [below=of d] {H};
    \path[->]
    (a) edge node {} (b)
    (a) edge node {} (c)
    (a) edge node {} (f)
    (b) edge node {} (e)
    (c) edge node {} (d)
    (d) edge node {} (h)
    (d) edge node {} (a)
    (e) edge node {} (f)
    (e) edge node {} (g)
    (f) edge node {} (b)
    (f) edge node {} (g)
    (h) edge node {} (g);
\end{tikzpicture}
\end{minipage}\hfill
\begin{minipage}[m]{0.5\linewidth}
Найти путь длины 13 из  \textbf{X} в \textbf{Y}
\end{minipage}
\end{frame}

\begin{frame}[noframenumbering]
  \transwipe[direction=90]
  \frametitle{Поиск путей в графах}
\begin{minipage}[m]{0.42\linewidth}
\begin{tikzpicture}[shorten >=1pt,node distance=2cm,on grid,auto]
   \node[state] (a) {\textbf{X}};
   \node[state] (b) [left=of a] {B};
   \node[state] (c) [right=of a] {C};
   \node[state] (d) [below=of c] {D};
   \node[state] (e) [below=of b] {E};
   \node[state] (f) [below=of a] {F};
   \node[state] (g) [below=of f] {\textbf{Y}};
   \node[state] (h) [below=of d] {H};
    \path[->]
    (a) edge node {} (b)
    (a) edge[red, line width=2pt] node {} (c)
    (a) edge[red, line width=2pt] node {} (f)
    (b) edge[red, line width=2pt] node {} (e)
    (c) edge[red, line width=2pt] node {} (d)
    (d) edge node {} (h)
    (d) edge[red, line width=2pt] node {} (a)
    (e) edge node {} (f)
    (e) edge[red, line width=2pt] node {} (g)
    (f) edge[red, line width=2pt] node {} (b)
    (f) edge node {} (g)
    (h) edge node {} (g);
\end{tikzpicture}
\end{minipage}\hfill
\begin{minipage}[m]{0.5\linewidth}
Найти путь длины 13 из  \textbf{X} в \textbf{Y}
\end{minipage}
\end{frame}

\begin{frame}[noframenumbering]
  \transwipe[direction=90]
  \frametitle{Поиск путей в графах}
\begin{minipage}[m]{0.42\linewidth}
\begin{tikzpicture}[shorten >=1pt,node distance=2cm,on grid,auto]
   \node[state] (a) {\textbf{X}};
   \node[state] (b) [left=of a] {B};
   \node[state] (c) [right=of a] {C};
   \node[state] (d) [below=of c] {D};
   \node[state] (e) [below=of b] {E};
   \node[state] (f) [below=of a] {F};
   \node[state] (g) [below=of f] {\textbf{Y}};
   \node[state] (h) [below=of d] {H};
    \path[->]
    (a) edge node {} (b)
    (a) edge node {} (c)
    (a) edge node {} (f)
    (b) edge node {} (e)
    (c) edge node {} (d)
    (d) edge node {} (h)
    (d) edge node {} (a)
    (e) edge node {} (f)
    (e) edge node {} (g)
    (f) edge node {} (b)
    (f) edge node {} (g)
    (h) edge node {} (g);
\end{tikzpicture}
\end{minipage}\hfill
\begin{minipage}[m]{0.5\linewidth}
    Найти кратчайший путь из  \textbf{X} в \textbf{Y}
\end{minipage}
\end{frame}

\begin{frame}[noframenumbering]
  \transwipe[direction=90]
  \frametitle{Поиск путей в графах}
\begin{minipage}[m]{0.42\linewidth}
\begin{tikzpicture}[shorten >=1pt,node distance=2cm,on grid,auto]
   \node[state] (a) {\textbf{X}};
   \node[state] (b) [left=of a] {B};
   \node[state] (c) [right=of a] {C};
   \node[state] (d) [below=of c] {D};
   \node[state] (e) [below=of b] {E};
   \node[state] (f) [below=of a] {F};
   \node[state] (g) [below=of f] {\textbf{Y}};
   \node[state] (h) [below=of d] {H};
    \path[->]
    (a) edge node {} (b)
    (a) edge node {} (c)
    (a) edge[red, line width=2pt] node {} (f)
    (b) edge node {} (e)
    (c) edge node {} (d)
    (d) edge node {} (h)
    (d) edge node {} (a)
    (e) edge node {} (f)
    (e) edge node {} (g)
    (f) edge node {} (b)
    (f) edge[red, line width=2pt] node {} (g)
    (h) edge node {} (g);
\end{tikzpicture}
\end{minipage}\hfill
\begin{minipage}[m]{0.5\linewidth}
Найти кратчайший путь из  \textbf{X} в \textbf{Y}
\end{minipage}
\end{frame}

\begin{frame}[fragile]
  \transwipe[direction=90]
  \frametitle{Поиск путей в графах: регулярные ограничения}
\begin{minipage}[m]{0.42\linewidth}
\begin{tikzpicture}[shorten >=1pt,node distance=2cm,on grid,auto]
   \node[state] (a) {\textbf{X}};
   \node[state] (b) [left=of a] {B};
   \node[state] (c) [right=of a] {C};
   \node[state] (d) [below=of c] {D};
   \node[state] (e) [below=of b] {E};
   \node[state] (f) [below=of a] {F};
   \node[state] (g) [below=of f] {\textbf{Y}};
   \node[state] (h) [below=of d] {H};
    \path[->]
    (a) edge node[above] {f} (b)
    (a) edge node {f} (c)
    (a) edge node {l} (f)
    (b) edge node[left] {t} (e)
    (c) edge node {o} (d)
    (d) edge node {o} (h)
    (d) edge node[above] {i} (a)
    (e) edge node {t} (f)
    (e) edge node[below] {w} (g)
    (f) edge node[above] {a} (b)
    (f) edge node {e} (g)
    (h) edge node {d} (g);
\end{tikzpicture}
\end{minipage}\hfill
\begin{minipage}[m]{0.5\linewidth}
\pause    Найти путь из \textbf{X} в \textbf{Y}, удовлетворяющий регулярному выражению $f \, o^*(d \mid l)$

\bigskip

\textcolor{white}{Constraints are a regular language}
\end{minipage}
\end{frame}

\begin{frame}[noframenumbering]
  \transwipe[direction=90]
  \frametitle{Поиск путей в графах: регулярные ограничения}
\begin{minipage}[m]{0.42\linewidth}
\begin{tikzpicture}[shorten >=1pt,node distance=2cm,on grid,auto]
   \node[state] (a) {\textbf{X}};
   \node[state] (b) [left=of a] {B};
   \node[state] (c) [right=of a] {C};
   \node[state] (d) [below=of c] {D};
   \node[state] (e) [below=of b] {E};
   \node[state] (f) [below=of a] {F};
   \node[state] (g) [below=of f] {\textbf{Y}};
   \node[state] (h) [below=of d] {H};
    \path[->]
    (a) edge node[above] {f} (b)
    (a) edge[red, line width=2pt] node {f} (c)
    (a) edge node {l} (f)
    (b) edge node[left] {t} (e)
    (c) edge[red, line width=2pt] node {o} (d)
    (d) edge[red, line width=2pt] node {o} (h)
    (d) edge node[above] {i} (a)
    (e) edge node {t} (f)
    (e) edge node[below] {w} (g)
    (f) edge node[above] {a} (b)
    (f) edge node {e} (g)
    (h) edge[red, line width=2pt] node {d} (g);
\end{tikzpicture}
\end{minipage}\hfill
\begin{minipage}[m]{0.5\linewidth}
    Найти путь из \textbf{X} в \textbf{Y}, удовлетворяющий регулярному выражению $f \, o^*(d \mid l)$
    \pause
\bigskip

Ограничения --- регулярный язык
\end{minipage}
\end{frame}

\begin{frame}[fragile]
  \transwipe[direction=90]
  \frametitle{Поиск путей в графах: КС ограничения}
\begin{minipage}[m]{0.42\linewidth}
\begin{tikzpicture}[shorten >=1pt,node distance=2cm,on grid,auto]
   \node[state] (a) {\textbf{X}};
   \node[state] (b) [left=of a] {B};
   \node[state] (c) [right=of a] {C};
   \node[state] (d) [below=of c] {D};
   \node[state] (e) [below=of b] {E};
   \node[state] (f) [below=of a] {F};
   \node[state] (g) [below=of f] {\textbf{Y}};
   \node[state] (h) [below=of d] {H};
   \path[->]
    (a) edge node[above] {(} (b)
    (a) edge node        {)} (c)
    (a) edge node        {(} (f)
    (b) edge node[left]  {(} (e)
    (c) edge node        {(} (d)
    (d) edge node        {(} (h)
    (d) edge node[above] {)} (a)
    (e) edge node        {)} (f)
    (e) edge node[below] {)} (g)
    (f) edge node[above] {)} (b)
    (f) edge node        {)} (g)
    (h) edge node        {)} (g);
\end{tikzpicture}
\end{minipage}\hfill
\begin{minipage}[m]{0.5\linewidth}
    Найти путь из \textbf{X} в \textbf{Y}, который является сбалансированной скобочной последовательностью

    \bigskip

    Ограничения --- контекстно-свободный язык, заданный грамматикой:

    \begin{center}
    $ S \rightarrow \varepsilon \mid ( S ) S $
    \end{center}

\bigskip

\begin{center}
\textcolor{white}{(())()} \textcolor{white}{()} \textcolor{white}{()()}
\end{center}

\end{minipage}
\end{frame}

\begin{frame}[noframenumbering]
  \transwipe[direction=90]
  \frametitle{Поиск путей в графах: КС ограничения}
\begin{minipage}[m]{0.42\linewidth}
\begin{tikzpicture}[shorten >=1pt,node distance=2cm,on grid,auto]
   \node[state] (a) {\textbf{X}};
   \node[state] (b) [left=of a] {B};
   \node[state] (c) [right=of a] {C};
   \node[state] (d) [below=of c] {D};
   \node[state] (e) [below=of b] {E};
   \node[state] (f) [below=of a] {F};
   \node[state] (g) [below=of f] {\textbf{Y}};
   \node[state] (h) [below=of d] {H};
   \path[->]
    (a) edge[red, line width=2pt] node[above] {(} (b)
    (a) edge node        {)} (c)
    (a) edge[green, line width=2pt] node        {(} (f)
    (b) edge[red, line width=2pt] node[left]  {(} (e)
    (c) edge node        {(} (d)
    (d) edge node        {(} (h)
    (d) edge node[above] {)} (a)
    (e) edge[red, line width=2pt] node        {)} (f)
    (e) edge[red, line width=2pt] node[below] {)} (g)
    (f) edge[red, line width=2pt] node[above] {)} (b)
    (f) edge[green, line width=2pt] node        {)} (g)
    (h) edge node        {)} (g);
\end{tikzpicture}
\end{minipage}\hfill
\begin{minipage}[m]{0.5\linewidth}
Найти путь из \textbf{X} в \textbf{Y}, который является сбалансированной скобочной последовательностью

\bigskip

Ограничения --- контекстно-свободный язык, заданный грамматикой:

\begin{center}
$ S \rightarrow \varepsilon \mid ( S ) S $
\end{center}

\bigskip

\begin{center}
\textcolor{red}{(())()} \hspace{1cm} \textcolor{green}{()} \hspace{1cm} ()()
\end{center}

\end{minipage}
\end{frame}


\begin{frame}[fragile]
  \transwipe[direction=90]
  \frametitle{Поиск путей в графах с КС ограничениями}
  Входные данные:

  \begin{itemize}
    \item ориентированный граф с метками на ребрах из алфавита $\Sigma$
    \item описание контекстно-свободного языка над алфавитом $\Sigma$
  \end{itemize}

\medskip

  Задача: найти все пути в графе, которые являются строками языка
\end{frame}

\begin{frame}[fragile]
  \transwipe[direction=90]
  \frametitle{Области применения}
\begin{itemize}
	\item Запросы к графовым базам данных
	\item Статический анализ кода
	\item Анализ семантических сетей
\end{itemize}
\end{frame}


\begin{frame}[fragile]
  \transwipe[direction=90]
  \frametitle{Синтаксический анализ для поиска путей}

Необходимо производить синтаксический анализ строк на путях

\bigskip

Обобщаем синтаксический анализ для работы с графами

\end{frame}


\begin{frame}
  \transwipe[direction=90]
  \frametitle{Существующие решения}

  \begin{itemize}
      \item Hellings Jelle ``Conjunctive context-free path queries'' 2014 --- теория КС запросов
      \item Xiaowang Zhang et al ``Context-free path queries on RDF graphs'' 2016 --- основано на CYK
      \item Kröni Daniel et al ``Parsing Graphs: Applying Parser Combinators to Graph Traversals'' 2013 --- комбинаторы для обходов графов
      \item Grigorev Semyon et al ``Context-free Path Querying with Structural Representation of Result'' 2017 --- основано на GLL
      \item Verbitskaia Ekaterina et al ``Relaxed parsing of regular approximations of string-embedded languages'' 2015 --- основано на RNGLR
  \end{itemize}

Не предоставляют средство прозрачной интеграции запросов в основной язык программирования

\end{frame}

\begin{frame}
  \transwipe[direction=90]
  \frametitle{Постановка задачи}

  Цель: разработка средства прозрачной интеграции КС запросов в язык программирования общего назначения

  Задачи:

  \begin{itemize}
      \item Разработать библиотеку парсер-комбинаторов для КС запросов, которая:
      \begin{itemize}
         \item поддерживает произвольные КС грамматики
         \item поддерживает произвольные входные графы
         \item не ограничивает формат входных данных
         \item позволяет в явном виде получать пути, удовлетворяющие ограничениям
         \item предоставляет механизм выполнения пользовательских семантических функций над путями
      \end{itemize}
      % \item Оценить применимость библиотеки для решения задач в прикладных областях
      \item Сравнить производительность реализованной библиотеки с существующими решениями
  \end{itemize}

\end{frame}

\begin{frame}[fragile]
  \transwipe[direction=90]
  \frametitle{Парсер-комбинаторы: прозрачная интеграция}
\begin{itemize}
  \item КС спецификация формулируется на языке программирования общего назначения
  \item Статические свойства основного языка затрудняют описание бессмысленных КС спецификаций
  \item Естественное представление результата
  \item Композициональность
\end{itemize}
\end{frame}

\begin{frame}[fragile]
  \transwipe[direction=90]
  \frametitle{Библиотека Meerkat}
\begin{itemize}
  \item Библиотека обобщенных парсер-комбинаторов на Scala
  \item Поддерживает произвольные КС спецификации (включая леворекурсивные и неоднозначные)
  \item Синтаксический анализ в стиле передачи продолжений с мемоизацией
  \item Вычисляет пользовательскую семантику
\end{itemize}

\begin{lstlisting}
val E: Nonterminal
= syn ( E ~ "^" ~ E
      | "-" ~ E
      | E ~ "*" ~ E
      | E ~ "/" ~ E
      | E ~ "+" ~ E
      | E ~ "-" ~ E
      | "(" ~ E ~ ")"
      | "[0-9]".r
      )
\end{lstlisting}
\end{frame}

\begin{frame}[fragile]
  \transwipe[direction=90]
  \frametitle{Поиск путей при помощи Meerkat}
\begin{minipage}[m]{0.42\linewidth}
\begin{tikzpicture}[shorten >=1pt,node distance=2cm,on grid,auto]
   \node[state] (a) {\textbf{X}};
   \node[state] (b) [left=of a] {B};
   \node[state] (c) [right=of a] {C};
   \node[state] (d) [below=of c] {D};
   \node[state] (e) [below=of b] {E};
   \node[state] (f) [below=of a] {F};
   \node[state] (g) [below=of f] {\textbf{Y}};
   \node[state] (h) [below=of d] {H};
   \path[->]
    (a) edge[red, line width=2pt] node[above] {(} (b)
    (a) edge node        {)} (c)
    (a) edge[green, line width=2pt] node        {(} (f)
    (b) edge[red, line width=2pt] node[left]  {(} (e)
    (c) edge node        {(} (d)
    (d) edge node        {(} (h)
    (d) edge node[above] {)} (a)
    (e) edge[red, line width=2pt] node        {)} (f)
    (e) edge[red, line width=2pt] node[below] {)} (g)
    (f) edge[red, line width=2pt] node[above] {)} (b)
    (f) edge[green, line width=2pt] node        {)} (g)
    (h) edge node        {)} (g);
\end{tikzpicture}
\end{minipage}\hfill
\begin{minipage}[m]{0.5\linewidth}
Найти путь из \textbf{X} в \textbf{Y}, который является сбалансированной скобочной последовательностью

    \bigskip

\begin{lstlisting}
val S: Nonterminal
= syn ( "(" ~ S ~ ")" ~ S
      | ""
      )
\end{lstlisting}

\bigskip

\begin{center}
\textcolor{red}{(())()} \hspace{1cm} \textcolor{green}{()} \hspace{1cm} ()()
\end{center}

\end{minipage}
\end{frame}


\begin{frame}[fragile]
  \transwipe[direction=90]
  \frametitle{Набор парсер-комбинаторов}

\begin{table}[h]
\small
\centering
\begin{tabular}{l@{}|l}
\multicolumn{1}{c|}{Комбинатор} & \multicolumn{1}{|c}{Описание} \\ \hline
{\lstinline!a ~ b!} & последовательное применение: {\lstinline!a!} затем {\lstinline!b!}   \\
{\lstinline!a | b!} & выбор: {\lstinline!a!} или {\lstinline!b!}         \\ \pause
{\lstinline!a ?!}   & опциональный разбор: {\lstinline!a!} или ничего   \\
{\lstinline!a *!}   & повторение, 0 или больше  {\lstinline!a!} \\
{\lstinline!a +!}   & повторение, 1 или больше {\lstinline!a!} \\
{\lstinline!a ^ f!} & применение функции {\lstinline!f!} к {\lstinline!a!}, если  {\lstinline!a!}~--- терминал \\
{\lstinline!a ^^!}  & вернуть результат применения {\lstinline!a!}, если  {\lstinline!a!}~--- терминал    \\
{\lstinline!a & f!} & применение функции {\lstinline!f!} к {\lstinline!a!}, если  {\lstinline!a!}~--- парсер \\
{\lstinline!a &&!}  &  вернуть результат {\lstinline!a!}, если  {\lstinline!a!}~--- парсер    \\
\hline
\end{tabular}
\end{table}
\end{frame}

\begin{frame}[fragile]
  \transwipe[direction=90]
  \frametitle{Реализация: входные данные}
\begin{center} Строка = линейный граф \end{center}
\begin{center}
\begin{tikzpicture}[shorten >=1pt,node distance=1.5cm,on grid,auto]
\tikzstyle{point} = [circle, minimum width = 3pt, fill]
   \node[point] (a) {};
   \node[point] (b)[right=of a] {};
   \node[point] (c)[right=of b] {};
   \node[point] (d)[right=of c] {};
   \node[point] (e)[right=of d] {};
   \node[point] (f)[right=of e] {};
   \node[point] (g)[right=of f] {};
   \path[->]
    (a) edge node {С} (b)
    (b) edge node {Т} (c)
    (c) edge node {Р} (d)
    (d) edge node {О} (e)
    (e) edge node {К} (f)
    (f) edge node {А} (g);
\end{tikzpicture}
\end{center}

\bigskip

\bigskip

\pause

\begin{center} В чем отличия графа от строки? \end{center}

\pause

\begin{multicols}{2}
\begin{center} Ветвления \end{center}

\begin{multicols}{2}
\begin{center}
\begin{tikzpicture}[shorten >=1pt,node distance=1.5cm,on grid,auto]
\tikzstyle{point} = [circle, minimum width = 3pt, fill]
   \node[point] (a) {};
   \node[point] (b)[right=of a, above right=of a] {};
   \node[point] (c)[right=of a, below=of b] {};
   \path[->]
    (a) edge node {A} (b)
    (a) edge[below] node {B} (c);
\end{tikzpicture}
\end{center}

\columnbreak

\begin{center}
\begin{tikzpicture}[shorten >=1pt,node distance=1.5cm,on grid,auto]
\tikzstyle{point} = [circle, minimum width = 3pt, fill]
   \node[point] (a) {};
   \node[point] (b)[left=of a, above left=of a] {};
   \node[point] (c)[left=of a, below=of b] {};
   \path[->]
    (b) edge node {A} (a)
    (c) edge[below] node {B} (a);
\end{tikzpicture}
\end{center}
\end{multicols}

\columnbreak

\pause

\begin{center} Циклы \end{center}
\begin{center}
\begin{tikzpicture}[shorten >=1pt,node distance=1.5cm,on grid,auto]
\tikzstyle{point} = [circle, minimum width = 3pt, fill]
   \node[point] (a) {};
   \node[point] (b)[right=of a] {};
   \node[point] (c)[below=of b] {};
   \node[point] (d)[left=of c] {};
   \path[->]
    (a) edge node {Y} (b)
    (b) edge node {A} (c)
    (c) edge node {D} (d)
    (d) edge node {A} (a);
\end{tikzpicture}
\end{center}
\end{multicols}
\end{frame}

\begin{frame}[fragile]
  \transwipe[direction=90]
  \frametitle{Общая идея}
\begin{itemize}
  \item Каждая позиция во входном потоке соответсвует набору промежуточных результатов синтаксического анализа
  \item В случае ветвления осуществляем синтаксический анализ независимо вдоль исходящих ребер
  \item Объединяем промежуточные результаты, если несколько путей приводят в одну вершину
  \item Не анализируем одни и те же позиции дважды
  \item Мемоизируем промежуточные результаты
\end{itemize}
\end{frame}


\begin{frame}[fragile]
  \transwipe[direction=90]
  \frametitle{Учет информации в вершинах}
\begin{minipage}[m]{0.42\linewidth}
\begin{tikzpicture}[shorten >=1pt,node distance=2cm,on grid,auto]
   \node[state] (a) {\textbf{X}};
   \node[state] (b) [left=of a] {B};
   \node[state] (c) [right=of a] {C};
   \node[state] (d) [below=of c] {D};
   \node[state] (e) [below=of b] {E};
   \node[state] (f) [below=of a] {F};
   \node[state] (g) [below=of f] {\textbf{Y}};
   \node[state] (h) [below=of d] {H};
   \path[->]
    (a) edge[red, line width=2pt] node[above] {(} (b)
    (a) edge node        {)} (c)
    (a) edge[green, line width=2pt] node        {(} (f)
    (b) edge[red, line width=2pt] node[left]  {(} (e)
    (c) edge node        {(} (d)
    (d) edge node        {(} (h)
    (d) edge node[above] {)} (a)
    (e) edge[red, line width=2pt] node        {)} (f)
    (e) edge[red, line width=2pt] node[below] {)} (g)
    (f) edge[red, line width=2pt] node[above] {)} (b)
    (f) edge[green, line width=2pt] node        {)} (g)
    (h) edge node        {)} (g);
\end{tikzpicture}
\end{minipage}\hfill
\begin{minipage}[m]{0.5\linewidth}
Найти путь из \textbf{X} в \textbf{Y}, который является сбалансированной скобочной последовательностью

    \bigskip

\begin{lstlisting}
val Query : Nonterminal
= syn (LV("X") ~ S ~
       LV("Y"))

val S: Nonterminal
= syn ( "(" ~ S ~ ")" ~ S
      | ""
      )
\end{lstlisting}

\bigskip

\begin{center}
\textcolor{red}{(())()} \hspace{1cm} \textcolor{green}{()} \hspace{1cm} ()()
\end{center}

\end{minipage}
\end{frame}

\begin{frame}[fragile]
  \transwipe[direction=90]
  \frametitle{Абстракция входных данных}
\begin{lstlisting}
trait Input[+E, +N] {
  def checkNode(nodeId: Int,
      predicate: N => Boolean): Option[N]
  def filterEdges(nodeId: Int,
      predicate: E => Boolean): Seq[(E, Int)] }
\end{lstlisting}

\bigskip

Мы предоставляем интеграцию с scalax.collection.Graph и Neo4j
\end{frame}

\begin{frame}[fragile]
  \transwipe[direction=90]
  \frametitle{Пример}
\begin{lstlisting}
val query =
  syn((
    syn(LV("Actor") ^^) ~ outLE("ACTS_IN") ~
    LV("Movie")) & (a => (a.name, a.toInt)))
executeQuery(query, input)
  .groupBy{ case(a, i) => i }
  .toIndexedSeq
  .map{ case(i, ms) => (ms.head._1, ms.length) }
  .sortBy{ case(a, mc) => -mc }}
  .take (10)
\end{lstlisting}
\end{frame}


\begin{frame}
  \transwipe[direction=90]
  \frametitle{Эксперименты}
\begin{itemize}
  \item Запросы к онтологиям
  \begin{itemize}
    \item Примерно в 3 раза быстрее, чем решение, основанное на GLL
  \end{itemize}
  \item Статический анализ кода: may-alias отношения
  \begin{itemize}
    \item До 5 раз быстрее, чем Trails
  \end{itemize}
  \item Запросы к базе данных ``Фильмы''
  \begin{itemize}
    \item До 10 раз медленее, чем Neo4j + Cypher
    \item Все запросы регулярны
  \end{itemize}
\end{itemize}

\end{frame}

\begin{frame}
  \transwipe[direction=90]
  \frametitle{Ограничения}
\begin{itemize}
  \item Накладные расходы при выполнении регулярных запросов
  \item Неизвестно, как вычислять произвольную семантику для найденных путей
  \begin{itemize}
    \item Если извлекать пути лениво, то в каком порядке?
    \item Как вычислять семантику, если в графе есть циклы?
  \end{itemize}
\end{itemize}
\end{frame}

\begin{frame}
  \transwipe[direction=90]
  \frametitle{Результаты}

  \begin{itemize}
      \item На языке программирования Scala реализована библиотека парсер-комбинаторов для запросов к графам
      % \item Показана применимость библиотеки для решения задач в прикладных областях
      \item Произведено сравнение производительности реализованной библиотеки с существующими решениями: она не уступает другим в выполнении контекстно-свободных запросов
      \item Основные результаты, полученные в данной работе были представлены на конференциях Scala Symposium 2018 (SCOPUS) и ``Языки программирования и компиляторы 2017'' (РИНЦ)
  \end{itemize}

\end{frame}

\end{document}
